\chapter{Introduction}

\section{Purpose}
This text is documentation for a software library for robot control called 
Simple Manipulator Control (SMC). 

The library has 2 purposes:
\begin{enumerate}
\item To speed up the development, testing and deployment
of (primarily) model-based controllers. This is facilitated via the 
``batteries included'' approach, 
by providing a ready-to-use collection of common control strategies.
\item To serve as an educational platform for those looking
	to get into robot control.
\end{enumerate}

These goals are achieved through the following.
The library is developed in Python. This choice allows
for short, flexible, easy-to-understand code, 
a wide array of available external packages,
while still providing real-time control, as heavy computation
is executed in C/C++ via pybind wrappers (via NumPy, Pinocchio etc).
Furthermore, the library is designed to be as simple as possible, 
both in terms of logic,
and the total amount of code. Its skeleton provides a common
interface to both real and simulated robots,
and a wrapper around control loops which are to be written by the users.
The library primarily targets manipulators, but also provides 
support for mobile manipulators,
mainly targeting whole-body control design and testing.

The simplicity comes at the cost of some programming comfort.
This is guided by a personal belief that robotics software 
should get out of the way
of the fact that a robotic manipulator is just 6+ motorized pendulums,
the control of which are easy to implement,
provided appropriate mathematical background.
Thus, the target audience for the library are control engineers
who want to prototype and test their control designs.
More broadly, the target audience  
are people with a more mathematical background
rather than a software background.

Regarding software.
Production-level robotics code should of course not be in Python,
or any language with garbage collection.
However, prototyping control should not have to be done
with the strictest possible software standards,
as this simply takes too much time, and it isn't even clear
that a given controller is worth implementing properly.
Having that said, poorly designed software is the 
biggest hindrance to any project including software.
Thus, the library aims to provide a sensible software architecture,
but without forcing the user to always stick to it.

To further facilitate the educational value
of the library, the documentation explains not only
the control design of existing controllers, but also
the specific software design principles and techniques used.
As mentioned above, software design is essential even for prototyping ---
any code longer than a few hundred lines requires high-level structure
to be useful.

\paragraph{Basic technical information}
\begin{itemize}
\item The library is implemented in Python.
\item The code is tested on MacOS, Arch Linux, Ubuntu 22 and 24, and should work on any Linux. Support for Windows is achieved
 through Docker. The installation procedure for any OS
is documented in the Installation section.
\item All of the robotics-related dependent libraries are written in C++ with Python bindings, proving fast execution despite the Python front-end
		(good enough for real-time control at >500Hz).
\item Most dependencies are optional, and are used only in
	dedicated modules.
\end{itemize}

% TODO: citations on different control strategies listed.
% prolly will mostly be books
\section{Overview of available robot types and controllers}

\paragraph{Available tobot types}
There is basic support for the following types of robots:
\begin{itemize}
		\item single arm manipulators, such as UR5e
		\item dual arm manipulators, such as YuMi
		\item single arm wheeled mobile manipulators
		\item dual arm wheeled mobile manipulators
\end{itemize}
\paragraph{Available controllers}
The library provides the implementation of following common 
velocity-level control strategies for point to
point motions, as well as trajectory tracking:
\begin{itemize}
\item joint space control
\item Cartesian space control via the following inverse kinematics solvers:
\begin{itemize}
	\item Jacobian-based solvers such as damped Jacobian pseudoinverse
	\item various QP formulations
\end{itemize} 
\item impedance control
\item optimal control in the form of trajectory optimization (including MPC)
\item dynamic motion primitives (DMP)
\end{itemize}
There is also some existing functionality for motion planning
\begin{itemize}
	\item path planning
\begin{itemize}
	\item dynamical system-based planning
	\item random path generation
\end{itemize} 
	\item interpolation
\begin{itemize}
	\item approximating B-splines
\end{itemize} 
	\item trajectory generation (ongoing work)
\begin{itemize}
	\item polynomial trajectories
\end{itemize} 
	\item visual servoing (ongoing work)
\end{itemize}

Future work consists of the following:
\begin{itemize}
	\item joint and Cartesian space torque-level control
\item controllers for extrinsic dexterity
\end{itemize}


\subsection{Supported robots}
\begin{itemize}
  \item UR robots
  \item ABB YuMi
  \item ABB mobile YuMi
  \item Sleipner + YuMi platform (SLuMi) 
  \item MiR mobile base with a UR arm (Heron)
\end{itemize}
Every robot, and its interface, has particular quirks
one has to be aware of when working with hardware.
To the extent pertinent for the library,
they are described in \ref{chapter-hardware}.

Planned extension of support:
\begin{itemize}
\item Franka Emika
\item KUKA LBR iiwa
\end{itemize}


\paragraph{Extending support to other robots}
Detailed instructions for extending to other robots are 
provided in Appendix \ref{sec-supportable-hardware}
but it consists of the following steps:
\begin{itemize}
\item Robot models are loaded via URDF.
\item The communication interface with the robot is achieved through either:
\begin{enumerate}
\item A direct interface, ex. socket with a protocol.
The interface has to contain functions for reading of the robot's state (joint positions, velocities, f/t sensor as vectors)
  and functions to send control inputs (joint velocities).
\item ROS1 driver. Due to ROS1 depending on deprecated 
	Python versions, the ROS1 driver is to run in a Docker
	translation layer, while communicating
	with SMC though protobuf-serialized messages over sockets.
\item ROS2 driver. The appropriate robot's ROS2 Node inherits
	SMC's RealRobotManager, providing both SMC and ROS2 functionality.
	In addition, in main files the controllers need to be registered
	instead of immediately executed,
	but otherwise everything is the same.
\end{enumerate}
\item Interfacing with grippers is dealt with on a case-by-case basis
	depending on the specific hardware setup and available drivers.
	Strong preference is given to independent gripper drivers
	as this allows for simultaneous control of the robot and
	the gripper.
\end{itemize}

\section{Usage instructions}
The best way to get a hang of the library is by looking at
the examples available in the \fbox{\texttt{examples}}
directory.
Different arguments and script description is available through 
\newline
\fbox{\texttt{python examples/script.py --help}}.
More involved examples come with a technical report explaining the 
goal of the example and how it is achieved.

The library itself is documented within this document.
This concerns the high-level design decisions,
and the most important elements.  
The \textbf{code itself is documented} and users are highly encouraged 
to look into it, as it can often be used as a starting point
when developing new functionality.

The available control designs are briefly introduced in the control chapter.
In those, the reader is referred to books and 
papers for their theoretical background.
Every controller comes with its example script,
and with an instruction of how to run it in
the \textbf{How to run} paragraphs, which denote where
the example is and which arguments are important for it.
The same is true for various robot classes.

A \textbf{suggested workflow} for writing new examples
and expanding the library can be found in the 
\ref{sec-suggested-workflow} section.

For users with little experience with Python, Linux and similar
topics there are references to relevant learning material 
in the Dependencies section. The idea is to provide 
the minimum working knowledge needed to write
a new controller and test it on hardware.
